% F081: precise mathematical teaching diagram.
\begin{tikzpicture}[x=1cm,y=1cm,>=stealth,font=\small]
\draw[rounded corners=3pt,gray!40] (0,0) rectangle (14.2,8.8);
\begin{scope}[yshift=1cm]
\node at (7.1,7.2) {也可以多次起伏：至少一个，不限制只有一个};
\begin{scope}[shift={(1.4,3.7)},x=2.45cm,y=1.5cm]
\draw[->,gray!65] (-0.2,0)--(4.45,0) node[right] {$x$};
\draw[->,gray!65] (0,-1.45)--(0,1.6) node[above] {$y$};
\draw[blue!75!black,line width=1.15pt,,domain=0:4,samples=160] plot (\x,{sin(180*\x)});
\fill[blue!75!black] (0,0) circle (1.8pt);\node[below left] at (0,0) {$A$};
\fill[blue!75!black] (4,0) circle (1.8pt);\node[below right] at (4,0) {$B$};
\draw[orange!85!black,line width=1.2pt] (0.28,1)--(0.72,1);
\fill[orange!85!black] (0.5,1) circle (1.8pt);
\draw[orange!85!black,line width=1.2pt] (1.28,-1)--(1.72,-1);
\fill[orange!85!black] (1.5,-1) circle (1.8pt);
\draw[orange!85!black,line width=1.2pt] (2.28,1)--(2.72,1);
\fill[orange!85!black] (2.5,1) circle (1.8pt);
\draw[orange!85!black,line width=1.2pt] (3.28,-1)--(3.72,-1);
\fill[orange!85!black] (3.5,-1) circle (1.8pt);
\end{scope}
\node[] at (7.1,1.2) {图例 $f(x)=\sin(\pi x),\quad 0\le x\le4$：有四处水平切线。};
\end{scope}
\node[text width=13cm,align=center] at (7.1,.9) {三种形状可以反复组合；一般曲线不必只有一段上升和一段下降。};
\end{tikzpicture}
